% !TeX root = ../drafts/c-flock-protocol.tex
\section{The Flock protocol}\label{annex:flock}

Flock~\cite{Flock26} proves many executions of one Boolean circuit at once. We use it for the eight instruction classes' circuits (\S\ref{sec:tables}, \S\ref{flock:classes}), each getting its own run of the protocol over its own witness. Since we do not follow their protocol exactly, we describe our version here, for one circuit of block size $\kblock$ with its constant wire at position $\jconst$; the eight differ in those two numbers alone, which \S\ref{flock:classes} gives, and the BLAKE2s compression, the largest at $\kblock=14$, is the running example.

\subsection{Statement and commitment}\label{flock:statement}

One instance of the circuit has a Boolean R1CS witness block of $2^{\kblock}$ bits. Its fixed matrices are $A_0,B_0,C_0\in\Ftwo^{2^{\kblock}\times 2^{\kblock}}$, with $C_0=I$. The witness includes the circuit inputs, outputs, AND-gate values, addition carries and one constant-one position; linear wires are substituted away. Within-block coordinates are the low-order coordinates and batch coordinates are the high-order coordinates throughout this annex, and a Boolean tuple is also read as an integer, lowest coordinate first.

For $2^{\kbatch}$ independent instances, the witness is:
\[
z:\cube{\kbool}\longrightarrow\Ftwo,\qquad \kbool=\kblock+\kbatch,
\]
and the batched matrices are
\[
A=I_{2^{\kbatch}}\otimes A_0,\qquad B=I_{2^{\kbatch}}\otimes B_0,\qquad C=I_{2^{\kbool}}
\]
Writing $a=Az$, $b=Bz$ and $c=Cz=z$, the R1CS statement is:
\begin{equation}\label{flock:eq:r1cs}
a(u)b(u)+c(u)=0\qquad\text{for every }u\in\cube{\kbool}
\end{equation}

We also enforce the position $\jconst$ of every block to be $1$ (preventing the all-zero witness):
\begin{equation}\label{flock:eq:const-one}
z(\jconst,\,t)=1\qquad\text{for every }t\in\cube{\kbatch}
\end{equation}

This second condition is enforced by lincheck (\S\ref{flock:lincheck}), and $\jconst$ is where the circuit puts its constant wire, the position after its ports (\S\ref{flock:classes}). Since it is asked of every block, the prover pads the batch out to $2^{\kbatch}$ with genuine witness blocks, the instances of its class's no-op on zero inputs, rather than with zeros.

Set $\kskip=6$ and $\nflock=\kbool-\kskip$. The prover commits to $z$ by packing its low $\kskip$ coordinates, $64$ bits, to a $\K$-element, following Ring-Switching (Annex~\ref{annex:rs}). The resulting multilinear $\qflock:\cube{\nflock}\to\K$ is committed as one region of the stack (\S\ref{sec:stacking}).


\subsection{Zerocheck}\label{flock:zerocheck}

Flock zerochecks \eqref{flock:eq:r1cs} (Definition~\ref{def:zerocheck}), but replaces the first $\kskip$ rounds by one univariate round~\cite{Gru24}. Split $u=(i,v)$ with $i\in\cube{\kskip}$ and $v\in\cube{\nflock}$. Let $\fembed:\F_{2^8}\hookrightarrow\E$ be the subfield embedding, with a byte interpreted in the polynomial basis, and define
\[
\skipdom=\{\fembed(0),\ldots,\fembed(63)\},\qquad \skipcoset=\fembed(64)+\skipdom=\{\fembed(64),\ldots,\fembed(127)\}.
\]
Thus $\skipdom$ is a $64$-point additive subspace and $\skipcoset$ is its disjoint coset, in the additive-NTT ordering of Annex~\ref{annex:pcs}.

Write $\lag_i$ for the Lagrange basis of $\skipdom$: $\lag_i(\fembed(i'))=[i=i']$.

\begin{definition}[Quirky extension]\label{def:quirky}
For $f:\cube{\kskip}\times\cube{n}\to\Ftwo$, its \emph{quirky extension} interpolates the first $\kskip$ coordinates over $\skipdom$ and extends the other $n$ multilinearly:
\[
  \qext f:\E\times\E^{n}\to\E,\qquad
  \qext f(Y,v):=\sum_{i=0}^{63}\lag_i(Y)\,\mle f(i,v),
\]
so $\qext f$ has degree $63$ in $Y$, degree one in every other variable, and $\qext f(\fembed(i),v)=\mle f(i,v)$ for every $i\in\cube{\kskip}$. Below $f$ is $a$, $b$ or $c$, with $n=\nflock$, or $A_0$ or $B_0$ read as a row followed by a column, with $n=2\kblock-\kskip$. Its argument $(Y,v)$ is a \emph{quirky point}.
\end{definition}

The zerocheck point $r \in \E^{\nflock}$ is defined as follows:
\begin{itemize}
  \item Its first seven coordinates $r_0, \dots, r_6$ are fixed to $\fembed(\mathtt{0xf7}),\fembed(\mathtt{0x53}),\fembed(\mathtt{0xb5})$ and $g_0^{2^j}/(1+g_0^{2^j})$ for $0\leq j<4$, where $g_0\in\E$ is public. Lemma~\ref{lem:fixed-zerocheck} applies.
  \item The remaining coordinates $r_7,\dots,r_{\nflock-1}$ are sampled at random by the verifier.
\end{itemize}

Define
\begin{equation}\label{flock:eq:pskip}
P(Y)=\sum_{v\in\cube{\nflock}}\eq(r,v)\bigl(\qext a(Y,v)\qext b(Y,v)+\qext c(Y,v)\bigr),
\end{equation}
which lies in $\E[Y]$ and has degree at most $126$, the product contributing $2\times 63$ and the linear term $63$. If \eqref{flock:eq:r1cs} holds, then $P(h)=0$ for every $h\in\skipdom$. The prover sends $P|_{\skipcoset}$, $64$ values: together with the $64$ assumed zeros on $\skipdom$ that is a degree-below-$128$ polynomial's worth of evaluations, so the verifier interpolates $P$ over the $128$-point space $\skipdom\cup\skipcoset$.

The verifier then samples $\zskip\in\E$ and reconstructs the claimed value
\[
v_P \qeq P(\zskip).
\]

The parties run sumcheck over the remaining $\nflock$ Boolean variables in \eqref{flock:eq:pskip}, with Gruen's eq-factor optimization~\cite{Gru24}. Writing $R_{\mathrm{zc}}\in\E$ for the running claim after the last round, the prover sends $v_a,v_b\in\E$ and the terminal identity
\begin{equation}\label{flock:eq:zc-terminal}
R_{\mathrm{zc}}=v_av_b+v_c
\end{equation}
is what \emph{defines} $v_c$: both sides solve it rather than transmit it. The zerocheck therefore ends leaving three claims at one point,
\begin{equation}\label{flock:eq:zc-claims}
\qext a(\zskip,\fc)=v_a,\qquad \qext b(\zskip,\fc)=v_b,\qquad \qext z(\zskip,\fc)=v_c,
\end{equation}
where $\fc\in\E^{\nflock}$ is the vector of sumcheck challenges and the third claim uses $c=z$.

\subsection{Lincheck}\label{flock:lincheck}

Lincheck reduces the three zerocheck claims to claims on $z$. The block structure is what keeps it cheap: the matrices repeat one circuit, so lincheck runs over a single block, not over the batch. Split $\fc=(\fc_{\mathrm{in}},\fc_{\mathrm{out}})$, where $\fc_{\mathrm{in}}\in\E^{\kblock-\kskip}$ selects the remaining within-block coordinates and $\fc_{\mathrm{out}}\in\E^{\kbatch}$ selects a batch instance.

\begin{lemma}[Block diagonality folds the batch into $z$]\label{lem:flock-block-diagonal}
For $a=Az$ with $A=I_{2^{\kbatch}}\otimes A_0$, at every point $(\zskip,\fc_{\mathrm{in}},\fc_{\mathrm{out}})\in\E\times\E^{\kblock-\kskip}\times\E^{\kbatch}$,
\[
\qext a(\zskip,\fc_{\mathrm{in}},\fc_{\mathrm{out}})=\sum_{j\in\cube{\kblock}}\qext A_0(\zskip,\fc_{\mathrm{in}},j)\,\mle z(j,\fc_{\mathrm{out}}),
\]
and likewise for $b=Bz$ with $B_0$ in place of $A_0$.
\end{lemma}

\begin{proof}
Block diagonality gives
\[
a(k,t)=\sum_{j\in\cube{\kblock}}A_0(k,j)\,z(j,t)\qquad\text{for }k\in\cube{\kblock},\;t\in\cube{\kbatch},
\]
so, splitting the row as $k=(\ksk,\kin)$ with $\ksk\in\cube{\kskip}$ and $\kin\in\cube{\kblock-\kskip}$ and expanding every extension by Fact~\ref{fact:eq},
\begin{align*}
\qext a(\zskip,\fc_{\mathrm{in}},\fc_{\mathrm{out}})&=\sum_{\ksk\in\cube{\kskip}}\lag_{\ksk}(\zskip)\,\mle a(\ksk,\fc_{\mathrm{in}},\fc_{\mathrm{out}})\\
&=\sum_{\ksk}\lag_{\ksk}(\zskip)\sum_{\kin}\eq(\fc_{\mathrm{in}},\kin)\sum_{t\in\cube{\kbatch}}\eq(\fc_{\mathrm{out}},t)\,a\bigl((\ksk,\kin),t\bigr)\\
&=\sum_{\ksk,\kin}\lag_{\ksk}(\zskip)\eq(\fc_{\mathrm{in}},\kin)\sum_{j\in\cube{\kblock}}A_0((\ksk,\kin),j)\underbrace{\sum_{t}\eq(\fc_{\mathrm{out}},t)\,z(j,t)}_{=\;\mle z(j,\fc_{\mathrm{out}})}\\
&=\sum_{j\in\cube{\kblock}}\Bigl(\underbrace{\sum_{\ksk,\kin}\lag_{\ksk}(\zskip)\,\eq(\fc_{\mathrm{in}},\kin)\,A_0((\ksk,\kin),j)}_{=\;\qext A_0(\zskip,\fc_{\mathrm{in}},j)}\Bigr)\mle z(j,\fc_{\mathrm{out}}).
\end{align*}
\end{proof}

The $c$ claim has the same shape, with the column index split as $j=(\jsk,\jin) \in\cube{\kskip} \times \cube{\kblock-\kskip}$, $\jsk$ the coordinates the univariate skip replaced and $\jin$ the within-block coordinates that remain:
\begin{equation}\label{flock:eq:c-as-column}
\qext z(\zskip,\fc_{\mathrm{in}},\fc_{\mathrm{out}})=\sum_{j\in\cube{\kblock}}\lag_{\jsk}(\zskip)\,\eq(\fc_{\mathrm{in}},\jin)\,\mle z(j,\fc_{\mathrm{out}}),
\end{equation}

\vspace{3mm}

The verifier samples $\lcalpha\in\E$ and batches the zerocheck's three claims \eqref{flock:eq:zc-claims} with the constant-position check \eqref{flock:eq:const-one}:
\begin{equation}\label{flock:eq:lincheck}
\begin{split}
v_a+\lcalpha v_b+\lcalpha^2v_c+\lcalpha^3\qeq\sum_{j\in\cube{\kblock}}\Bigl(&\qext A_0(\zskip,\fc_{\mathrm{in}},j)+\lcalpha\qext B_0(\zskip,\fc_{\mathrm{in}},j)\\
&+\lcalpha^2\lag_{\jsk}(\zskip)\eq(\fc_{\mathrm{in}},\jin)+\lcalpha^3[j=\jconst]\Bigr)\mle z(j,\fc_{\mathrm{out}}).
\end{split}
\end{equation}

\vspace{2mm}

The parties start by running $\kblock - \kskip = 14 - 6 = 8$ sumcheck rounds on $\eqref{flock:eq:lincheck}$, binding the $\jin$ coordinates (leaving the skip index $\jsk$ unfolded), with challenges $\fc'_{\mathrm{in}}$ in coordinate order (the reverse of their round order). Let $R_{\mathrm{lc}}\in\E$ be the resulting sumcheck claim.

The prover sends $(s_0,\ldots,s_{63})\in\E^{64}$, claiming $s_i=\mle z(i,\fc'_{\mathrm{in}},\fc_{\mathrm{out}})$. These 64 claims will later be proven by Ring-Switching.

The verifier finally needs to check:

\begin{equation}\label{flock:eq:lincheck-terminal}
\begin{split}
R_{\mathrm{lc}}\qeq\sum_{i=0}^{63}\Bigl(&\qext A_0(\zskip,\fc_{\mathrm{in}},i,\fc'_{\mathrm{in}})+\lcalpha\,\qext B_0(\zskip,\fc_{\mathrm{in}},i,\fc'_{\mathrm{in}})\\
&+\lcalpha^2\lag_i(\zskip)\,\eq(\fc_{\mathrm{in}},\fc'_{\mathrm{in}})+\lcalpha^3[i=0]\,\eq(\fc'_{\mathrm{in}},8)\Bigr)s_i.
\end{split}
\end{equation}

\subsubsection{Concluding Lincheck}\label{flock:native}


With the row split $k=(\ksk,\kin) \in\cube{\kskip} \times \cube{\kblock-\kskip}$ of the proof above,
\begin{equation}\label{flock:eq:erow}
\qext A_0(\zskip,\fc_{\mathrm{in}},j)=\sum_{k\in\cube{\kblock}}\erow(k)\,A_0(k,j),\qquad
\erow(\ksk,\kin)=\lag_{\ksk}(\zskip)\,\eq(\fc_{\mathrm{in}},\kin),
\end{equation}
and likewise for $B_0$.

Writing $\wcol(\jsk,\jin)=s_{\jsk}\,\eq(\fc'_{\mathrm{in}},\jin)$ with $(\jsk,\jin) \in\cube{\kskip} \times \cube{\kblock-\kskip}$, each matrix term in $\eqref{flock:eq:lincheck-terminal}$ is simplified as follows:
\begin{equation}\label{flock:eq:matrix-form}
\sum_{i=0}^{63}\qext A_0(\zskip,\fc_{\mathrm{in}},i,\fc'_{\mathrm{in}})\,s_i=\sum_{k,j\in\cube{\kblock}}A_0(k,j)\,\erow(k)\,\wcol(j).
\end{equation}

One forward walk of the circuit returns \eqref{flock:eq:matrix-form} for both matrices (\S\ref{flock:matrices}), never touching the millions of nonzero entries they hold: one multiplication per committed wire on the $A$ side, one per $\wedge$ wire on the $B$ side, and one for the constant factor the remaining $B$ rows share, in $O(\text{circuit})$ work. For the BLAKE2s compression the committed wires are $864$ free input bits (over $896$ positions, the $32$-bit finalization word leaving half of its word empty), the constant, $256$ committed $\oplus$ wires carrying the result, and $14{,}878$ $\wedge$ wires: one per carry of the $480$ $32$-bit additions of the ten rounds, less the two that the first round's constant lanes make structurally zero (\S\ref{flock:classes}). The $c$ term of \eqref{flock:eq:lincheck-terminal} is $\lcalpha^2\inner{\erow}{\wcol}$ in the same notation, and needs no walk: both sides are tensors, so it contracts to $\lcalpha^2\eq(\fc_{\mathrm{in}},\fc'_{\mathrm{in}})\sum_i\lag_i(\zskip)s_i$.

\subsubsection{Ring switching}\label{flock:ringswitch}

The role of Ring-Switching is now to prove:
\[
s_i \qeq \mle z(i,\fc'_{\mathrm{in}},\fc_{\mathrm{out}})\qquad\text{for }i\in\cube{\kskip}.
\]


\subsection{Protocol summary}\label{flock:summary}

\begin{protocol}[Flock summary]\label{flock:protocol}
After the witness $z$ has been committed, and its size announced:
\begin{enumerate}
\item The verifier forms the zerocheck point $r$: seven fixed coordinates, the rest sampled. The $\kskip$ variables the univariate skip replaced consume no equality challenge.
\item The prover sends $P|_{\skipcoset}$, $64$ values. The verifier samples $\zskip$ and reconstructs $v_P$.
\item The parties run $\nflock$ quadratic zerocheck rounds on \eqref{flock:eq:pskip}, which the prover ends by sending $v_a,v_b$; the terminal identity \eqref{flock:eq:zc-terminal} then fixes $v_c$.
\item The verifier samples $\lcalpha$ and the parties run the $\kblock-\kskip$ quadratic lincheck rounds on \eqref{flock:eq:lincheck}. The prover sends $(s_0,\ldots,s_{63})$, and the verifier checks \eqref{flock:eq:lincheck-terminal} with one circuit walk.
\end{enumerate}
What Flock leaves is one family of $64$ claims for Ring-Switching (\S\ref{flock:ringswitch}): $(s_i \qeq \mle z(i, \fc'_{\mathrm{in}},\fc_{\mathrm{out}}))_{0 \leq i \leq 63}$.
\end{protocol}

\subsection{Evaluating the matrices}\label{flock:matrices}

$A_0$ and $B_0$ encode a Boolean circuit, but are never materialized.

\subsubsection{The circuit}\label{flock:circuit}

\begin{definition}[Circuit]\label{def:circuit}
A \emph{circuit} is an ordered sequence of wires $w_1,\ldots,w_N$. Each wire is either a \emph{source}, a $\oplus$ wire with a possibly empty set $S_w$ of earlier wires, or a $\wedge$ wire with an ordered pair $(u_w,u'_w)$ of earlier wires. Its value lies in $\Ftwo$: a $\oplus$ wire is the sum of the values in $S_w$, and a $\wedge$ wire is the product of the values of $u_w$ and $u'_w$. The source $w_1$ is the constant $1$ and the other sources are free inputs.

A set of \emph{committed} wires containing all sources, all $\wedge$ wires, and a subset of the $\oplus$ wires (typically the result of the computation) is equipped with an injection $w\mapsto\posof{w}$ into $\cube{\kblock}$. The constant has position $\posof{w_1}=\jconst$. A position $k\in\cube{\kblock}$ is \emph{empty} if $k\neq\posof{w}$ for every committed wire $w$. A witness block $z_t=z(\cdot,t)$ stores each committed wire's value at its position and stores $0$ at empty positions.
\end{definition}

For $p\in\cube{\kblock}$, write $[\,\cdot=p\,]:\cube{\kblock}\to\Ftwo$ for the function that equals $1$ at $p$ and $0$ at every other position. The \emph{expansion} of a wire over committed positions is
\begin{equation}\label{flock:eq:expand}
\lform{w}=[\,\cdot=\posof{w}\,]\ \text{ if $w$ is committed},\qquad \lform{w}=\sum_{u\in S_w}\lform{u}\ \text{ otherwise},
\end{equation}
where the second case applies only to uncommitted $\oplus$ wires. By induction, $\inner{\lform{w}}{z_t}$ is the value of $w$.

\subsubsection{The R1CS matrices}\label{flock:rows}

For every committed wire $w$, assign two sets of wires to the corresponding row:
\[
(S^A_w,S^B_w)=
\begin{cases}
(\{u_w\},\{u'_w\}) & \text{if $w$ is a $\wedge$ wire},\\
(S_w,\{w_1\}) & \text{if $w$ is a $\oplus$ wire},\\
(\{w\},\{w_1\}) & \text{if $w$ is a source}.
\end{cases}
\]
and define:
\begin{equation}\label{flock:eq:row-lists}
A_0(\posof{w},\cdot)=\sum_{u\in S^A_w}\lform{u},\qquad B_0(\posof{w},\cdot)=\sum_{u\in S^B_w}\lform{u}.
\end{equation}
For every empty position $k$, set $A_0(k,\cdot)=B_0(k,\cdot)=0$. This specifies every row of both matrices.

\vspace{3mm}

The R1CS constraints given by $A_0$, $B_0$ and $C_0=I$ faithfully encode the circuit.
\subsubsection{Forward algorithm}\label{flock:forward}

Given a column vector $W:\cube{\kblock}\to\E$, the forward algorithm computes the two row-indexed vectors $A_0W$ and $B_0W$ without constructing either matrix. The native lincheck verifier uses $W=\wcol$ and retains only the scalar $\inner{\erow}{A_0W}+\lcalpha\inner{\erow}{B_0W}$.

To compute them, visit the wires from $w_1$ to $w_N$ and evaluate every expansion against $W$:
\[
\inner{\lform{w}}{W}=W(\posof{w})\ \text{ if $w$ is committed},\qquad \inner{\lform{w}}{W}=\sum_{u\in S_w}\inner{\lform{u}}{W}\ \text{ otherwise};
\]
then \eqref{flock:eq:row-lists} gives
\begin{equation}\label{flock:eq:row-pair}
\inner{A_0(\posof{w},\cdot)}{W}=\sum_{u\in S^A_w}\inner{\lform{u}}{W},\qquad \inner{B_0(\posof{w},\cdot)}{W}=\sum_{u\in S^B_w}\inner{\lform{u}}{W}.
\end{equation}
Together with zero at every empty row, these values are the desired vectors.

\subsubsection{Transpose algorithm}\label{flock:backward}

For $X\in\{A,B\}$, the transpose algorithm computes the complete column-indexed vector $\colmar^{X}=X_0^{\!\top}\erow$ without constructing $X_0$. The prover needs all $2^{\kblock}$ entries of both vectors:
\begin{equation}\label{flock:eq:marginal}
\colmar^{X}(j)=\qext X_0(\zskip,\fc_{\mathrm{in}},j)=\sum_{k\in\cube{\kblock}}X_0(k,j)\erow(k)=\bigl(X_0^{\!\top}\erow\bigr)(j),\qquad X\in\{A,B\},
\end{equation}
where the second equality uses \eqref{flock:eq:erow}. Viewing each row as a column-indexed vector, empty rows contribute zero, so \eqref{flock:eq:row-lists} gives
\begin{align*}
\colmar^{X}
&=\sum_{w\ \mathrm{committed}}\erow(\posof{w})\sum_{u\in S^X_w}\lform{u}\\
&=\sum_u\left(\sum_{\substack{w\ \mathrm{committed}\\u\in S^X_w}}\erow(\posof{w})\right)\lform{u}.
\end{align*}
The last equality exchanges the finite sums and collects every occurrence of the same expansion $\lform{u}$. Thus, for every circuit wire $u$, define
\begin{equation*}
\chg{u}^{X}=\sum_{\substack{w\ \mathrm{committed}\\u\in S^X_w}}\erow(\posof{w}),
\qquad \text{so that}\qquad
\colmar^{X}=\sum_u\chg{u}^{X}\lform{u}.
\end{equation*}
Run the following procedure independently for $X=A$ and $X=B$. Initialize the output table $\colmar^{X}$ to zero and visit $w_N,\ldots,w_1$. At wire $w$, eliminate the term $\chg{w}^{X}\lform{w}$ using the appropriate case of \eqref{flock:eq:expand}:
\begin{itemize}
\item if $w$ is committed, then $\chg{w}^{X}\lform{w}=\chg{w}^{X}[\,\cdot=\posof{w}\,]$, so add $\chg{w}^{X}$ to $\colmar^{X}(\posof{w})$;
\item otherwise, $w$ is an uncommitted $\oplus$ wire and
\[
\chg{w}^{X}\lform{w}=\chg{w}^{X}\sum_{u\in S_w}\lform{u}=\sum_{u\in S_w}\chg{w}^{X}\lform{u},
\]
so add $\chg{w}^{X}$ to the existing coefficient $\chg{u}^{X}$ of every predecessor $u\in S_w$.
\end{itemize}

\subsection{The class circuits}\label{flock:classes}

Each instruction class's circuit (\S\ref{sec:tables}) is a circuit in the sense of Definition~\ref{def:circuit}, built from a small vocabulary of gates over $64$-bit words: XOR, AND, the multiplexer $s\,(x\oplus y)\oplus y$, the ripple-carry adder and the carry-save multiplier below. One block layout serves all eight: the ports first, $64$ positions per word, the inputs then the outputs, bit $i$ of a word at position $i$ of its range; the constant wire at $\jconst$, the position right after them; then the $\wedge$ wires in the order the circuit makes them. A port's bits are free wires when it is an input, and committed $\oplus$ wires when it is an output, so that the result leaves the circuit; a port bit with no wire, such as the $63$ high bits of a one-bit output, is an empty position, which \eqref{flock:eq:r1cs} forces to zero. Each port is therefore one packed $\K$-element, equal to the register, bytecode or memory word it stands for (\S\ref{sec:tab-class}). A hint (\texttt{DIV}'s quotient and remainder) is an input port that no tuple carries.
\begin{center}
\begin{tabular}{lrrrrr}
\hline
Circuit & ports (in $\to$ out) & $\kblock$ & $\jconst$ & $\wedge$ wires & positions used\\
\hline
\texttt{ALU}   & $4\to2$  & $10$ & $384$   & $424$      & $809$\\
\texttt{LOAD}  & $4\to2$  & $10$ & $384$   & $320$      & $705$\\
\texttt{STORE} & $5\to3$  & $10$ & $512$   & $346$      & $859$\\
\texttt{SHIFT} & $4\to1$  & $10$ & $320$   & $579$      & $900$\\
\texttt{MUL}   & $3\to1$  & $12$ & $256$   & $2{,}175$  & $2{,}432$\\
\texttt{MULH}  & $3\to1$  & $13$ & $256$   & $4{,}545$  & $4{,}802$\\
\texttt{DIV}   & $5\to2$  & $13$ & $448$   & $5{,}156$  & $5{,}605$\\
\texttt{HASH}  & $14\to4$ & $14$ & $1{,}152$ & $14{,}878$ & $16{,}031$\\
\hline
\end{tabular}
\end{center}
What the two verifiers and the prover have to agree on is this layout and the order the $\wedge$ wires are made in, which fixes every position; the order of the $\oplus$ wires is free, since they are substituted away.

\paragraph{Addition.} A ripple-carry adder over $n$ bits ($n=64$, or $32$ in the compression). With $c_0=0$, the carry into position $i+1$ is the majority of $a_i,b_i,c_i$,
\[
c_{i+1}=(a_i\oplus c_i)(b_i\oplus c_i)\oplus c_i ,
\]
one $\wedge$ wire, and the sum bit $a_i\oplus b_i\oplus c_i$ is affine in the wires before it. The carry out of the top position falls off the modulus, which leaves $n-1$ products. A subtraction is the addition of the complement with a carry in of $1$, and a comparison is its carry out.

\paragraph{Multiplication.} A schoolbook multiplier would pay one product per partial product $a_ib_j$ before summing them. Over $\Ftwo$ the partial products are free: with $e_{ij}=1\oplus a_i\oplus b_j$, one has $2a_ib_j=a_i+b_j-1+e_{ij}$ over the integers. Summing over $i,j$, with $\bar a=2^{64}-1-a$ the complement,
\[
2ab=\sum_{i<64}r_i\,2^{i}+\bar a+\bar b+(a+b)\,2^{64}+1-2^{128},\qquad r_i=\begin{cases}b&a_i=1\\ \bar b&a_i=0,\end{cases}
\]
whose right side is a sum of rows of bits each affine in the inputs. Its column $0$ is $e_{00}+\bar a_0+\bar b_0+1=2+2g$ with $g=\bar a_0\bar b_0$, so after that one product the identity halves: $ab$ is $1+g$ plus the other columns shifted down one place, which is $66$ rows, the $64$ rows $r_i$ and one each for $a$ and $b$, with $1$ and $g$ in the empty low positions of two of them.

Carry-save steps then compress the rows. A step replaces three rows by their sum row, the position-wise $\oplus$, and their carry row, the position-wise majority shifted up one place: a majority costs one product where at least two of the three rows have a bit, except that where exactly two do and the carry row is still free at that position, one of the two bits moves into it for nothing. Each step takes the three rows that end lowest. After $64$ steps two rows remain, and the adder above finishes. The low word of the product keeps $64$ positions of every row and drops the carry out of the top one, which is \texttt{MUL}'s $2{,}175$ products; the high word keeps all $128$, which is \texttt{MULH}'s, its signed forms correcting the product of the operands' magnitudes. \texttt{DIV} multiplies the divisor by the hinted quotient the same way and adds the remainder.

\paragraph{The compression.} \texttt{HASH} is the BLAKE2s compression itself: the ten rounds' $80$ G functions, each six $32$-bit additions on the halves of the ports' words, a three-operand addition being two of them chained, and XORs and rotations, which are free. Only the carries are products, $31$ per addition.

\paragraph{Batch size.} The zerocheck of \S\ref{flock:zerocheck} fixes seven coordinates of its point beyond the $\kskip$ skipped ones, so it needs $\kbool\ge13$, and lincheck works on batches of at least eight blocks, $\kbatch\ge3$. For the circuits of $\kblock\ge10$ the second bound is the binding one; a batch of the $\kblock=10$ circuits therefore has at least eight instances, and every table at least eight rows (\S\ref{sec:tables}). Lincheck runs $\kblock-\kskip$ rounds, from $4$ to $8$.

\subsection{Soundness analysis}\label{flock:soundness}


The commitment determines $\qflock$, hence $z$, before any challenge is drawn, and the PCS proves the claims opened on $\qflock$. So take the one transmitted family to be true of $z$,
\[
s_i=\mle z(i,\fc'_{\mathrm{in}},\fc_{\mathrm{out}}),\qquad i\in\cube{\kskip},
\]
up to the PCS error and the $2^{-160}$ of Annex~\ref{annex:rs}, and let us show a $z$ failing \eqref{flock:eq:r1cs} or \eqref{flock:eq:const-one} is rejected.

Lincheck comes first, and it is the only thing that checks anything: the zerocheck's own terminal identity defines $v_c$ instead of testing it, so every lie upstream reaches lincheck as a wrong triple $(v_a,v_b,v_c)$. Its terminal check \eqref{flock:eq:lincheck-terminal} is evaluated from the true $s$ and the matrices themselves, by the walk of \S\ref{flock:native}, so a false \eqref{flock:eq:lincheck} survives the $\kblock-\kskip$ rounds with probability at most $2(\kblock-\kskip)/|\E|$ (Fact~\ref{fact:sumcheck}, degree two in each round). And by Lemma~\ref{lem:flock-block-diagonal} and \eqref{flock:eq:c-as-column}, \eqref{flock:eq:lincheck} says
\[
\bigl(v_a+\qext a(\zskip,\fc)\bigr)+\lcalpha\bigl(v_b+\qext b(\zskip,\fc)\bigr)+\lcalpha^2\bigl(v_c+\qext z(\zskip,\fc)\bigr)+\lcalpha^3\bigl(1+\mle z(\jconst,\fc_{\mathrm{out}})\bigr)=0,
\]
degree three in $\lcalpha$, all four coefficients fixed before $\lcalpha$ is drawn. Acceptance therefore pins them all, except with $3/|\E|$ (Lemma~\ref{lem:sz}): $v_a$, $v_b$ and $v_c$ are the true evaluations, and $\mle z(\jconst,\fc_{\mathrm{out}})=1$, which is \eqref{flock:eq:const-one} except with $\kbatch/|\E|$ more, $\fc_{\mathrm{out}}$ being random.

With all three pinned, the zerocheck's terminal identity \eqref{flock:eq:zc-terminal} becomes a real check again: it fixes $R_{\mathrm{zc}}$ to the true value of the summand at $\fc$.

Now let $z$ violate \eqref{flock:eq:r1cs} at $u=(i,v)$. Then $P(\fembed(i))$ is the extension at $r$ of a nonzero Boolean multilinear (Definition~\ref{def:quirky}, Fact~\ref{fact:eq}), and the seven fixed coordinates of $r$ carry $\Ftwo$-independent weights, so it stays nonzero except with $(\nflock-7)/|\E|$ (Lemma~\ref{lem:fixed-zerocheck}): the $64$ zeros the verifier assumes on $\skipdom$ are false. Let $\hat P$ interpolate those zeros together with the received coset values. It vanishes on $\skipdom$ while $P$ does not, so $\hat P\neq P$ and $v_P\neq P(\zskip)$ except with $127/|\E|$ (Corollary~\ref{cor:idtest}). The $\nflock$ zerocheck rounds then run on a false initial claim with a true terminal value, and reject except with $2\nflock/|\E|$.

Summing those terms, such a $z$ is accepted with probability at most
\[
\frac{2(\kblock-\kskip)+3+\kbatch+(\nflock-7)+127+2\nflock}{|\E|}=\frac{5\kblock+4\kbatch+93}{|\E|}<2^{-183},
\]
on top of the two errors granted above.
